In the remaining part of this paper we show several examples of polynomial games and recovered their equilibria by the method explained in the \cref{sec:poly}. We implemented the hierarchy of programs \eqref{momsdp} in Julia using \href{https://jump.dev/}{JuMP} \cite{JuMP} and the \href{https://jump.dev/SumOfSquares.jl/stable/}{SumOfSquares.jl} package \cite{PolyJuMP}. We solved the semidefinite programs using \href{https://www.mosek.com/}{Mosek} on a laptop with a 3.1GHz dual-core CPU. Our code is available on GitHub at \href{https://github.com/votroto/PolyNets.jl}{https://github.com/votroto/PolyNets.jl}.

\begin{example}[Two-player zero-sum polynomial game]
    Consider a two-player zero-sum game with strategy sets $S_i$ given by
    \[
        S_{i} = \{ x \in \R \mid  -x^2 + 1 \ge 0 \}.
    \] 
    The utility functions are
    \[
        u_{1}(x_{1},x_{2}) = 2x_{1}x_{2}^2 - x_{1}^2 - x_{2} = -u_{2}(x_{1},x_{2}).
    \]
    Since polynomial network games are a generalization of polynomial games, we can use the same framework to solve them. A pure NE is achieved when Player 1 plays $x_{1} = 0.4$ and Player~2 plays $x_{2} = 0.63$. The optimal payoffs to each player are $\bw = (-0.47, 0.47)$. This problem was solved at order 4 of the hierarchy. Due to the low number of variables the problem required only 44 constraints and it was solved within 10\,ms.
\end{example}

\begin{example}[Three-player game with a complete graph]
    We consider a polynomial network game with the player set $N = \{1, 2, 3\}$. Each strategy set $S_i$ is equal to
    \[
        S_{i}   = \{ x \in \R \mid  -x^2 + 1 \ge 0 \}
    \]
    and the utility functions are
    \begin{equation*}\begin{aligned}
        u_{1}(x_{1},x_{2},x_{3}) &= -2x_{1}x_{2}^2 - 2x_{1}^2 + 5x_{1}x_{2} - 4x_{1}x_{3} - x_{2} - 2x_{3},\\
        u_{2}(x_{1},x_{2},x_{3}) &= 2x_{1}x_{2}^2 - 2x_{2}x_{3}^2 - 2x_{1}^2 - 5x_{1}x_{2} - 2x_{2}^2 + 5x_{2}x_{3} + x_{2},\\
        u_{3}(x_{1},x_{2},x_{3}) &= 2x_{2}x_{3}^2 + 4x_{1}^2 + 4x_{1}x_{3} + 2x_{2}^2 - 5x_{2}x_{3} + 2x_{3}.\\
    \end{aligned}\end{equation*}
    A partially mixed Nash equilibrium is this strategy profile:
        \[\begin{aligned}
        x_1 &= -0.06\\
        x_2 &= 0.35\\
        x_3 &=\begin{cases}
            1 &72\%\\
            -1 &28\%\\
        \end{cases}
    \end{aligned}\]
    The corresponding utilities are $\bw = (-1.22, 0.26, 0.97)$. This problem converged at order 4, resulting in 96 constraints. The solve time was around 10\,ms.
\end{example}

\begin{example}[Polynomial security game]
   We consider a polynomial security game with attackers $N_A=\{1,2\}$, defenders $N_D=\{3,4\}$, and targets $T_3=\{t_1,t_2\}$ and $T_4=\{t_3,t_4\}$. The capacities of attackers are $a_1=0.5$, $a_2=0.8$, and the defenders' capacities are $d_3=d_4=1$. The efficiency functions are:
    \begin{equation*}\begin{aligned}
        e_{t1}(x,y) &= (1-y^2)\cdot 0.7x^2\\
        e_{t2}(x,y) &= (y^2 -2y +1)\cdot (1.4x - 0.7x^2)\\
        e_{t3}(x,y) &= (1-y^2)\cdot (1.8x - 0.9x^2)\\
        e_{t4}(x,y) &= (y^2 -2y +1)\cdot 0.9x^2\\
    \end{aligned}\end{equation*}
    Our algorithm returns the partially mixed strategy profile at order 6:
    \[\begin{aligned}
        (x_{1,t_1}, x_{1,t_2}, x_{1,t_3}, x_{1,t_4}) &=
            (0.29, 0.0, 0.0, 0.21) \\
        (x_{2,t_1}, x_{2,t_2}, x_{2,t_3}, x_{2,t_4}) &=
            (0.36, 0.13, 0.04, 0.24) \\
        (y_{1,t_1}, y_{1,t_2}) &=\begin{cases}
            (0.33, 0.67) &46\%\\
            (0.69, 0.31) &54\%\\
        \end{cases}\\
        (y_{2,t_2}, y_{2,t_4}) &=\begin{cases}
            (0.46, 0.54) &2\%\\
            (0.08, 0.92) &6\%\\
            (0.91, 0.09) &92\%\\
        \end{cases}
    \end{aligned}\]
    The payoffs to each player are $\bw = (0.1, 0.3, -0.17, -0.23)$. While the payoffs converged at order 4 already and took only 3 seconds to solve, the corresponding strategies could only be extracted at order 6, which took 6 minutes to solve. Due to the high dimensionality of the utility functions and the fast growth of the hierarchy, the resulting SDP contained over 25 thousand constraints.
\end{example}

While running the experiments, we noticed the payoffs computed by this method are generally of high quality regardless of the solver used; The corresponding measures, on the other hand, were less accurate. For example, the strategy in the last example yields a loss of only $0.16$ instead of $0.17$. This appears to be a problem inherent to our approach rather than something that could be fine-tuned within the solvers.
